\begin{theorem}[Realization of one terminal Algorithm 3.1 step]
Let \(G=(V,E)\) be a finite loopless multigraph, let \(e_0\in E\),
and let \(\mathcal S\) be a certified Tashkinov series with \(r\)
noninitial iterations, palette size \(k\), input coloring \(\varphi\),
and data \(\mathcal D\).  Put
\[
 i=\operatorname{last}(r)\in\operatorname{Fin}(r+1),\qquad
 T=\mathcal D.\operatorname{tree}(\operatorname{last}(r+1)).
\]
Suppose that \(\Delta(G)+1\leq k\), that
\(\mathcal D.\operatorname{coloring}(i)\) is proper and leaves exactly
\(e_0\) uncolored, and that \(V(T)\) is closed but not strongly closed
under this coloring.

Then there exist a tree-sequence \(U\), a partial coloring \(\psi\),
finite sets \(S\subseteq\operatorname{Fin}(k)\) and \(F\subseteq E\),
an optional extension kind \(\Theta\), optional vertices \(v,u\),
optional colors \(\delta,\gamma\), and a certified Tashkinov series
\(\mathcal S^+\) with \(r+1\) noninitial iterations such that
\[
 \mathcal S^+.\operatorname{data}
   =\operatorname{TashkinovSeriesDataAppend}
      (\mathcal D,U,\psi,S,F,\Theta,v,u,\delta,\gamma)
 \quad\text{and}\quad
 \operatorname{length}(T)<\operatorname{length}(U).
\]
Thus every old tree, coloring, stored set, connecting-edge set, kind,
vertex, and color slot of \(\mathcal D\) occurs in \(\mathcal S^+\) at
its canonical old index, while the displayed values are the exact new
terminal slots.
\end{theorem}
\begin{proof}
The successor of \(i=\operatorname{last}(r)\) is the last tree index of
the old data, so the tree used at coloring index \(i\) is exactly \(T\).
Apply the choice trichotomy \noderef{AlgorithmThreeOneChoiceDataExists}
at this index.  Its hypotheses are precisely the degree bound,
properness, unique-uncolored-edge condition, closure, and failure of
strong closure assumed here.  We realize its three alternatives in turn.

First suppose that it returns a revisiting certificate \(R\).  The fields
of \noderef{RevisitingExtensionCertificate} give an exact edge identity
\(f=R.\operatorname{edge}\) in \(\partial_G(V(T))\), its inside and
outside endpoints, and the two colors \(R.\delta,R.\gamma\).  By
\noderef{TreeSequenceBoundaryEdgeAppendExists}, there is a tree-sequence
\(A\) such that
\[
             \operatorname{TreeSequenceAdjoinEdge}(T,f,A).       \tag{1}
\]
In particular, by \noderef{TreeSequenceAdjoinEdge}, \(A\) has length
\(\operatorname{length}(T)+1\), contains \(T\) as the exact ordered
initial segment, and has the same exact edge identity \(f\) as its new
last edge.  Apply \noderef{TashkinovClosureExists} to obtain a closure
\(U\) of \(A\) under the old terminal coloring.  By
\noderef{TashkinovClosure}, there is a Tashkinov augmentation sequence
from \(A\) to \(U\); hence
\noderef{TashkinovAugmentationSequenceInitialSegment} says that \(A\)
is an exact ordered initial segment of \(U\).

We spell out the resulting composition because the final step constructor
records it.  The length inequalities for \(T\leq A\) and \(A\leq U\)
compose.  At every old vertex index the entry of \(U\) equals that of
\(A\), which equals that of \(T\); the same is true at every old edge
index.  These are exactly the clauses of
\noderef{TreeSequenceInitialSegment}, so \(T\leq U\).  Moreover,
\[
 \operatorname{length}(T)+1=\operatorname{length}(A)
       \leq\operatorname{length}(U),
\]
and therefore \(\operatorname{length}(T)<\operatorname{length}(U)\).

For this branch choose
\[
 \psi=\mathcal D.\operatorname{coloring}(i),\quad
 S=\{R.\delta,R.\gamma\},\quad F=\{f\},\quad
 \Theta=\mathrm{RE},
\]
and choose the optional vertices and colors to be respectively
\(R\)'s inside endpoint, outside endpoint, delta, and gamma.  Let
\(\mathcal D^+\) be the exact append supplied by
\noderef{TashkinovSeriesDataAppend}.  The old index \(i\) embeds as the
predecessor of the new last index.  The backward-history data used there
are unchanged: \noderef{TashkinovStoredColorsAppend} identifies its
computed stored-color set with the old one, and
\noderef{RevisitingExtensionCertificateAppend} transports \(R\) to a
certificate over \(\mathcal D^+\), retaining the exact colors, edge, and
endpoints.  The preceding hypotheses, (1), the closure of \(A\) to
\(U\), the composed initial-segment relation, and the strict inequality
therefore fill, field for field, the RE constructor of
\noderef{TashkinovSeriesStep} at the new last index.  Applying
\noderef{TashkinovSeriesAppend} gives a certified series
\(\mathcal S^+\) whose data are literally \(\mathcal D^+\).  This proves
the conclusion in the RE case.

Suppose next that the trichotomy returns the SE alternative.  Write its
maximum-defective choice as
\[
                Q=(\pi,v,f,u,\delta).
\]
By the fields of \noderef{MaximumDefectiveChoice} and the definition
\noderef{DefectiveEdge}, the exact selected edge \(f\) belongs to a
boundary-color set and hence to \(\partial_G(V(T))\).  Apply
\noderef{TreeSequenceBoundaryEdgeAppendExists} to obtain an exact
one-edge append \(A=T+f\), and apply
\noderef{TashkinovClosureExists} to obtain a closure \(U\) of \(A\)
under \(\pi\).  Exactly the composition argument in the RE case, using
\noderef{TashkinovAugmentationSequenceInitialSegment}, proves
\(T\leq U\) and
\[
                     \operatorname{length}(T)<\operatorname{length}(U).
                                                                    \tag{2}
\]

Choose
\[
 \psi=\pi,\quad S=\{\delta\},\quad F=\{f\},\quad
 \Theta=\mathrm{SE},\quad
 (v,u,\delta,\gamma)=(v,u,\delta,\mathrm{none}),
\]
where the last tuple denotes the four optional entries.  Form the exact
append \(\mathcal D^+\).  The equality in
\noderef{TashkinovStoredColorsAppend} transports the maximum choice,
including its selected stability and maximality clauses, to the embedded
old slots of \(\mathcal D^+\).  It also transports the universal SE
disjointness alternative returned by the trichotomy.  If an RE
certificate existed at that embedded index, the reverse direction of
\noderef{RevisitingExtensionCertificateAppend} would give one for the
old data, contradicting the SE alternative's no-RE clause.  Thus the
old properness, unique-uncoloredness, closure and non-strong-closure
hypotheses, the transported choice and SE predicate, the exact append
\(A=T+f\), its closure \(U\), the order relation, and (2) give all fields
of the SE constructor of \noderef{TashkinovSeriesStep}.  A final
application of \noderef{TashkinovSeriesAppend} gives the required exact
data equality and strict growth.

It remains to realize the PE alternative.  Retain its no-RE proof, its
witnessed failure of the SE predicate, the same choice
\(Q=(\pi,v,f,u,\delta)\), the color \(\gamma\) satisfying
\noderef{ParallelExtensionColor}, and the proper coloring \(\sigma\)
which leaves exactly \(e_0\) uncolored, is stable for the enlarged
controlled set, and satisfies
\[
                    \operatorname{TreeExit}(T,\sigma,\gamma,\delta,v).
                                                                    \tag{3}
\]

We first extract a second exact defective edge.  The selected-edge field
of \(Q\), expanded through \noderef{DefectiveEdge}, supplies some color
whose \noderef{BoundaryColorEdges} set contains \(f\) and is defective.
Since \(\pi(f)=\delta\), equality of the two colored values on this same
exact edge identity identifies that color with \(\delta\).  Hence
\(\delta\) is a \noderef{DefectiveColor} and its exact boundary-color
set has at least two elements.  Choose
\[
 g\in\partial_{\pi,\delta}(V(T)),\qquad g\ne f.              \tag{4}
\]
This is a choice of a distinct edge identity; no endpoint-pair
identification is made.

Every boundary edge is incident with a vertex of \(T\) by
\noderef{BoundaryEdges}.  Since \(\delta\) belongs to the inserted part
of the controlled color set, the edge clause of
\noderef{StableColoring} fixes both \(f\) and \(g\) between \(\pi\) and
\(\sigma\).  Consequently
\[
                       \sigma(f)=\sigma(g)=\delta.             \tag{5}
\]
The missing-color clause of \noderef{ParallelExtensionColor} says that
\(\gamma\) is missing at \(v\) under \(\pi\), and vertexwise stability
makes it missing there under \(\sigma\).  If \(\gamma=\delta\), then
(5) and the incidence of \(f\) with \(v\) would put \(\gamma\) in the
\noderef{PresentColors} at \(v\), contradicting its membership in the
complementary set \noderef{MissingColors}.  Thus
\[
                              \gamma\ne\delta.                 \tag{6}
\]

Let
\[
 K=\operatorname{TwoColorComponent}(G,\sigma,\gamma,\delta,v).
\]
We claim that the edge-node of \(g\) does not belong to \(K\).  If it
did, use (4) and \noderef{BoundaryEdges} to orient the endpoint pair of
\(g\) with an inside endpoint \(w\in V(T)\).  By (5), the incidence
edge between the vertex-node \(w\) and edge-node \(g\) belongs to the
\((\gamma,\delta)\)-incidence graph of
\noderef{TwoColorIncidenceGraph}.  Hence \(w\) also belongs to the same
component \(K\).  Apply
\noderef{TreeExitComponentMeetsTreeOnlyAtExit} to (3), the properness of
\(\sigma\), and the fact that \(\gamma\) is missing at \(v\).  It gives
\(w=v\).  Now the distinct exact edges \(g\) and \(f\) are adjacent at
\(v\) and both have color \(\delta\) by (5), contradicting
\noderef{IsProperEdgeColoring}.  This proves
\[
                              g\notin K.                        \tag{7}
\]

Define the new coloring by the component swap
\[
             \psi=\operatorname{TwoColorComponentSwap}
                    (G,\sigma,\gamma,\delta,v).
\]
By (6), \noderef{TwoColorComponentSwapProper} makes \(\psi\) proper.
The definition \noderef{TwoColorComponentSwap} also preserves the
exact-one-uncolored property directly.  Indeed, the old uncolored edge
\(e_0\) still has value \(\mathrm{none}\), whether or not its edge-node
is tested against \(K\).  For any \(a\ne e_0\), choose its old colored
value \(c\) using \noderef{IsOneEdgeUncolored}.  Outside \(K\) its new
value is still \(c\); inside \(K\) it is \(\delta\) when
\(c=\gamma\), is \(\gamma\) when \(c=\delta\), and is still \(c\)
otherwise.  In every case it remains colored.  Moreover, (7) and (5)
give
\[
                              \psi(g)=\delta.                   \tag{8}
\]

We next prove that \(\delta\) is missing at \(v\) under \(\psi\).
Suppose instead that an edge identity \(a\) incident with \(v\) has
\(\psi(a)=\delta\).  If its edge-node is outside \(K\), then the swap
fixes it, so \(\sigma(a)=\delta\).  Its incidence with the root
vertex-node \(v\) then puts it in \(K\), a contradiction.  If its
edge-node is in \(K\), inspection of the swap and (6) shows that the only
old value which can become \(\delta\) is \(\gamma\).  Thus
\(\sigma(a)=\gamma\), contradicting that \(\gamma\) is missing at
\(v\).  No such \(a\) exists, and the complementary definitions
\noderef{PresentColors} and \noderef{MissingColors} give
\[
                    \delta\in\overline\psi(v).                 \tag{9}
\]

The vertex \(v\) belongs to \(V(T)\) by the defective-vertex field of
\noderef{MaximumDefectiveChoice}.  Equations (4), (8), and (9) therefore
contradict the defining condition \noderef{ColoredVertexSetClosed} for
\(V(T)\) under \(\psi\).  Hence \(T\) is not closed under \(\psi\).
Apply \noderef{TashkinovClosureExists} directly to \(T\), obtaining a
closure \(U\) under \(\psi\).  Its augmentation sequence makes
\(T\leq U\) by
\noderef{TashkinovAugmentationSequenceInitialSegment}.  This sequence
cannot be empty, because then \(U=T\) would be closed, contrary to the
preceding paragraph.  A nonempty augmentation sequence has a last
\noderef{TashkinovAugmentation}, which increases length by one after an
initial segment of length at least that of \(T\).  Consequently
\[
                       \operatorname{length}(T)<\operatorname{length}(U).
                                                                    \tag{10}
\]

For PE choose
\[
 S=\{\delta,\gamma\},\quad F=\{f\},\quad
 \Theta=\mathrm{PE},\quad
 (v,u,\delta,\gamma)=(v,u,\delta,\gamma)
\]
as the optional terminal entries, and form the exact append
\(\mathcal D^+\) with new coloring \(\psi\) and new tree \(U\).
As in SE, \noderef{TashkinovStoredColorsAppend} transports the maximum
choice, its witnessed failure of the universal SE predicate, and the
stability relation to the embedded old slots.  The no-RE clause transports
backward by \noderef{RevisitingExtensionCertificateAppend}.  The
history-sensitive choice of \(\gamma\) transports exactly by
\noderef{ParallelExtensionColorAppend}.  Thus (3), properness and
unique-uncoloredness of \(\sigma\), the definition and the two proved
properties of \(\psi\), the closure of \(T\) to \(U\), its ordered
initial-segment certificate, and (10) supply every field of the PE
constructor of \noderef{TashkinovSeriesStep}.  Applying
\noderef{TashkinovSeriesAppend} produces the required certified series
with data exactly \(\mathcal D^+\).

In this last construction \(f\) is recorded exactly in the new singleton
connecting-edge set, but the closure was taken directly from \(T\); no
claim that \(f\) belongs to \(U\) has been made.  The three trichotomy
alternatives are exhaustive, so the theorem follows.
\end{proof}
